\chapter{System design and architecture}

\section{Application structure}

Figure~\ref{fig:arch} shows the structure. A browser talks to one Go process over HTTP. The process serves two interfaces from the same services: a JSON API under \code{/api} and server-rendered HTML pages. PostgreSQL stores all state. The optional mail server is the only other external dependency, and e-mail can be switched to a log-only mode for demonstration.

\begin{figure}[H]
\centering
\begin{tikzpicture}[
  box/.style={draw,rounded corners=2pt,fill=black!4,align=center,minimum height=1.15cm,inner sep=5pt,font=\small},
  arr/.style={-{Stealth[length=2.2mm]},thick},
  node distance=1.0cm and 1.05cm]
\node[box] (browser) {\textbf{Browser}\\ HTML pages, htmx\\ JSON clients};
\node[box,right=of browser] (api) {\textbf{Go application}\\ chi router, middleware\\ handlers $\to$ services};
\node[box,right=of api] (pg) {\textbf{PostgreSQL 16}\\ tables, constraints,\\ views};
\node[box,below=0.9cm of api] (mail) {\textbf{SMTP server}\\ optional e-mail copies};
\draw[arr] (browser) -- (api);
\draw[arr] (api) -- (pg);
\draw[arr] (api) -- (mail);
\end{tikzpicture}
\caption{Application structure and responsibility boundaries}
\label{fig:arch}
\end{figure}

The code follows a layered architecture, shown in Figure~\ref{fig:layers}. Each request passes through middleware (security headers, request logging, rate limiting for login, authentication and role checks), a handler that parses the request and formats the response, a service that contains the business rules, and a repository that contains the SQL. Dependencies point downwards only. Services depend on repository \emph{interfaces} declared next to them, so the same services run on PostgreSQL repositories in production and on in-memory repositories in tests.

\begin{figure}[H]
\centering
\begin{tikzpicture}[
  layer/.style={draw,rounded corners=2pt,fill=black!4,minimum width=8.6cm,minimum height=0.85cm,align=center,font=\small},
  arr/.style={-{Stealth[length=2mm]},thick},
  node distance=0.42cm]
\node[layer] (mw) {Middleware: authentication, \code{RequireRole}, CSRF origin check, rate limit};
\node[layer,below=of mw] (h) {Handlers (JSON API) and web pages (HTML templates)};
\node[layer,below=of h] (s) {Services: business rules, validation, \code{AccessControl}};
\node[layer,below=of s,fill=black!10] (r) {Repository interfaces\\[-1pt]{\scriptsize implemented by PostgreSQL repositories and in-memory repositories}};
\node[layer,below=of r] (db) {PostgreSQL (constraints, indexes, views)};
\foreach \a/\b in {mw/h,h/s,s/r,r/db} \draw[arr] (\a) -- (\b);
\end{tikzpicture}
\caption{Layers of the application; dependencies point downwards}
\label{fig:layers}
\end{figure}

The main source directories are \code{internal/handlers} and \code{internal/web} (presentation), \code{internal/services} (rules), \code{internal/repository} and \code{internal/memrepo} (persistence), \code{internal/models} (shared types) and \code{internal/app}, which assembles the router so that the production binary and the tests use exactly the same routes. The commands \code{cmd/api}, \code{cmd/seed} and \code{cmd/loadtest} provide the server, a demonstration-data loader and a load generator.

\section{Responsibility boundaries}

Handlers contain no business rules. They translate a request into a service call and a service error into a status code through one helper: a forbidden action becomes 403, a missing object 404 and everything else that is a validation error 400. Services receive an explicit \code{Actor} value holding the user identifier and role instead of reading the HTTP request. This makes the access rules testable by constructing an actor in a unit test.

\subsection*{Centralised access control}
\label{sec:access}

The component \code{AccessControl} answers two questions used by attendance, grades and the timetable: may this actor record data for this class, and may this actor read the records of this pupil? An administrator may do both. A teacher may record for a class if they are its homeroom teacher or have at least one timetable lesson with it, and may read the pupils of such classes. A pupil may read only their own records. A parent may read the records of pupils linked to them. Everything else is denied. Because every service calls the same two functions, adding a feature cannot silently introduce a different interpretation of ``the teacher of a class''.

\subsection*{Server-rendered interface}

The interface is generated by Go \code{html/template} and embedded into the binary with \code{go:embed}, so the container needs no separate static-file server and the build has no JavaScript toolchain. The htmx library is vendored and served from the application itself, which allows a strict content-security policy that permits scripts only from the same origin. Forms are ordinary HTML forms that work without scripting; htmx only replaces the page body to avoid a full reload.

\section{Relational data model}

The schema has fourteen tables, three reporting views and five migrations. Figure~\ref{fig:er} shows the principal entities and Table~\ref{tab:models} their purpose and constraints.

\begin{figure}[H]
\centering
\begin{tikzpicture}[
  ent/.style={draw,rounded corners=2pt,fill=black!4,align=center,minimum width=2.5cm,minimum height=0.95cm,font=\small,inner sep=4pt},
  rel/.style={-{Stealth[length=2mm]},thick},
  lab/.style={font=\scriptsize,midway,fill=white,inner sep=1pt}]
\node[ent] (users) {\textbf{users}\\ role, e-mail, hash};
\node[ent,right=2.2cm of users] (classes) {\textbf{classes}\\ homeroom teacher};
\node[ent,below=1.5cm of users] (students) {\textbf{students}\\ class, number};
\node[ent,right=2.2cm of students] (timetable) {\textbf{timetable}\\ day, time, room};
\node[ent,below=1.5cm of students] (grades) {\textbf{grades}\\ mark 1--5, term};
\node[ent,right=2.2cm of grades] (attendance) {\textbf{attendance}\\ status per day};
\node[ent,left=2.0cm of students] (news) {\textbf{news}\\ title, body, pinned};
\draw[rel] (students) -- (users) node[lab,pos=.5] {is};
\draw[rel] (students) -- (classes) node[lab,pos=.45] {in};
\draw[rel] (timetable) -- (classes) node[lab,pos=.5] {for};
\draw[rel] (grades) -- (students) node[lab,pos=.5] {of};
\draw[rel] (attendance) -- (students) node[lab,pos=.4] {of};
\draw[rel] (news) -- (users) node[lab,pos=.5] {author};
\end{tikzpicture}
\caption{Principal entities (teachers, parents, subjects, link tables and notifications omitted)}
\label{fig:er}
\end{figure}

\begin{table}[H]
\centering
\caption{Principal tables and their constraints}
\label{tab:models}
\small
\begin{tabularx}{\textwidth}{@{}l>{\raggedright\arraybackslash}X>{\raggedright\arraybackslash}p{4.6cm}@{}}
\toprule
Table & Purpose & Important constraint \\
\midrule
users & Account, role, active flag, soft-delete timestamp & Unique e-mail; role is an enumerated type \\
teachers, students, parents & Role-specific data referring to a user & Unique \code{user\_id}; foreign key to users \\
classes & Class name, level, section, homeroom teacher & Foreign key to users \\
subjects & Subject name and code & Unique code \\
teacher\_subjects, student\_parents & Qualifications and parent--child links & Composite keys \\
timetable & Weekly lesson of a class & Check: end after start; day 0--6 \\
attendance & One status per pupil, class and day & Unique (pupil, class, date) \\
grades & A mark with subject, class, term and author & \code{chk\_grades\_five\_point}: mark in $\{1,\dots,5\}$ \\
notifications & Message to a user with read flag & Index on unread messages \\
news & School news item & Index for the feed (pinned, then newest) \\
\bottomrule
\end{tabularx}
\end{table}

The schema is in third normal form. Account data (e-mail, password hash, name, role) is stored once in \code{users}; the role tables keep only role-specific columns. Names of classes, subjects and teachers are never copied into marks, attendance or timetable rows, they are reached through foreign keys. Averages and percentages are not stored but computed, so they cannot become stale. One deliberate denormalisation exists: \code{grades} and \code{attendance} record the \code{class\_id} the pupil belonged to \emph{at the time}, because a pupil may later change class and the historical fact must not move with them. Users are soft-deleted so that the marks of a removed account are preserved.

\section{Integrity and query design}

Primary keys, foreign keys, unique and check constraints enforce different aspects of consistency \cite{pgconstraints}. The five-point rule is an example of defence in depth: the service rejects an invalid mark with a message that a teacher can act on, and the check constraint ensures that no other code path, script or manual SQL statement can store one.

Indexes follow the dominant queries: marks by pupil and term for the diary and reports, timetable by teacher and day for conflict detection, notifications by user and read state for the bell counter, attendance by pupil and date, and the news feed ordered by pinned flag and creation time. Three views (\code{v\_student\_\allowbreak grade\_\allowbreak summary}, \code{v\_student\_\allowbreak attendance\_\allowbreak summary}, \code{v\_class\_\allowbreak overview}) provide aggregated marks and attendance for reporting tools without duplicating the logic of Chapter~2 outside the application; the attendance view uses the same formula as Equation~(3).

\section{Security design}

\textbf{Authentication.} Passwords are hashed with bcrypt at its default cost. A successful login returns a signed JSON Web Token valid for 24 hours; the web interface stores it in an \code{HttpOnly}, \code{SameSite=Strict} cookie (with the \code{Secure} flag behind HTTPS). On every web request the user is loaded again, so a deactivated or deleted user is signed out immediately even if the token has not expired.

\textbf{Cross-site request forgery.} Every state-changing web request (anything other than GET and HEAD) must carry an \code{Origin} or \code{Referer} header whose host equals the request host; otherwise it is rejected with 403. This complements the \code{SameSite} cookie attribute \cite{owaspcsrf}.

\textbf{Injection and scripting.} SQL statements are parameterised. Templates escape all dynamic text, and the news body is stored and displayed as plain text split into escaped paragraphs, so an administrator cannot accidentally publish markup. Responses carry a content-security policy (\code{default-src 'self'}, \code{script-src 'self'}, \code{frame-ancestors 'none'}), \code{X-Frame-Options: DENY}, \code{X-Content-Type-Options: nosniff} and a same-origin referrer policy \cite{mdncsp}.

\textbf{Brute force.} Login requests are limited per client address with a fixed window (10 per minute by default, configurable); excess requests receive 429 with a \code{Retry-After} header. Public self-registration does not exist.

\section{Persistence, scheduling and deployment}

\textbf{Notifications and reminders.} Recording a mark or an absence creates in-app notifications for the pupil and the linked parents. An optional mailer sends a copy by SMTP or, in demonstration mode, writes it to the log. A reminder job runs at a configurable hour (15:00 by default) and notifies teachers who have not recorded attendance for the day.

\textbf{Migrations.} Schema changes are plain SQL files in \code{migrations/}, each with an \code{up} and a \code{down} script in the format of the migration tool \cite{migrate}. The continuous-integration pipeline applies all scripts forwards, backwards and forwards again against a clean database.

\textbf{Container.} A multi-stage Docker build compiles the server and the seeding tool with the Go toolchain and copies the binaries into a small Alpine image that runs as a non-root user and includes time-zone data, so that the reminder job follows local school time. Docker Compose starts the application together with PostgreSQL 16.

\textbf{Backup and restore.} The script \code{scripts/backup.sh} writes a compressed custom-format dump with \code{pg\_dump} \cite{pgdump} to a temporary name, renames it when complete, verifies it with \code{pg\_restore --list} and removes dumps older than a configurable number of days. The script \code{scripts/restore.sh} restores a dump into an empty database and stops at the first error. The objectives are a recovery point of 24 hours (daily dump) and a recovery time of under 30 minutes for a school-sized database; Section~\ref{sec:backup} reports a restore test.
